\chapter{Introduction}

Carpooling has emerged as a sustainable solution to urban transportation challenges, including traffic congestion, air pollution, and rising fuel costs. Traditional carpooling platforms rely on centralized architectures, where a single authority controls user data, transactions, and trust mechanisms. This centralization introduces several vulnerabilities including data privacy concerns, single points of failure, lack of transparency in pricing and operations, and dependence on intermediaries for dispute resolution.

Blockchain technology offers a paradigm shift by enabling decentralized, transparent, and secure systems. By integrating blockchain with carpooling services, we can create a trustless environment where users interact directly through smart contracts, eliminating the need for intermediaries while ensuring data integrity and transparency.

D-CARPOOL is designed as a comprehensive solution that addresses these challenges by combining the best of centralized efficiency with decentralized security and transparency. The platform implements seven distinct blockchain security mechanisms to provide multi-layered protection for users.

\section{Motivation}

The motivation for developing D-CARPOOL stems from several key observations:

\begin{itemize}
    \item \textbf{Campus Mobility Needs}: IIIT Kottayam students require efficient transportation solutions for daily commutes, airport transfers, and intercity travel. Traditional ride-sharing apps often have limited coverage in smaller cities.
    
    \item \textbf{Trust Issues}: Existing platforms lack transparent reputation systems, leading to safety concerns and uncertain user experiences. There is no mechanism to verify user credentials without exposing sensitive information.
    
    \item \textbf{Cost Efficiency}: Centralized platforms charge significant commissions, increasing the cost for users. A decentralized approach can reduce these costs substantially.
    
    \item \textbf{Data Privacy}: Users are increasingly concerned about how their personal data is used and stored by centralized platforms. Sensitive information like location, travel patterns, and personal contacts need protection.
    
    \item \textbf{Emergency Safety}: Traditional platforms lack robust emergency mechanisms with distributed trust. Single-point-of-failure emergency systems are vulnerable to manipulation.
    
    \item \textbf{Blockchain Potential}: The immutability and transparency of blockchain technology, combined with advanced cryptographic protocols, can create a more trustworthy environment for peer-to-peer transactions.
\end{itemize}

\section{Problem Statement}

Traditional carpooling systems face several critical challenges:

\begin{enumerate}
    \item \textbf{Security Concerns}: Passengers lack detailed information about drivers, creating safety risks. The absence of verifiable credentials and transparent history increases vulnerability. There is no privacy-preserving way to verify driver licenses.
    
    \item \textbf{Centralized Control}: Single point of failure in centralized systems can lead to service disruptions. Platform operators have complete control over user data and policies.
    
    \item \textbf{Lack of Transparency}: Hidden algorithms for pricing and matching reduce user trust. Dispute resolution processes are opaque and biased toward the platform.
    
    \item \textbf{Limited Accountability}: Reputation systems can be manipulated without immutable records. No permanent record of user behavior across platforms.
    
    \item \textbf{Privacy Violations}: Sensitive data like vehicle information is visible to all users. Emergency contact data is stored in plain text, vulnerable to breaches.
    
    \item \textbf{No Audit Trail}: Ride history can be modified or deleted, making dispute resolution difficult. No efficient mechanism to prove historical rides occurred.
\end{enumerate}

\section{Objectives}

The primary objectives of this project are:

\begin{itemize}
    \item To develop a secure, decentralized carpooling platform for IIIT Kottayam students
    \item To implement blockchain-based identity management using Soulbound Tokens (SBT)
    \item To create Zero-Knowledge Proof verification for driver credentials
    \item To establish on-chain hash integrity for ride agreements
    \item To implement DPoS governance for transparent dispute resolution
    \item To develop Merkle Tree audit trails for efficient ride history verification
    \item To integrate CP-ABE encryption for fine-grained access control
    \item To deploy Shamir's Secret Sharing for emergency data protection
    \item To provide a user-friendly web interface for seamless interaction
    \item To reduce operational costs by eliminating intermediaries
    \item To implement comprehensive emergency SOS features
\end{itemize}

\section{Scope}

The scope of D-CARPOOL encompasses:

\subsection{Current Implementation}

\begin{itemize}
    \item Email-based authentication with OTP verification restricted to institutional emails
    \item Ride creation, search, and management functionalities
    \item Real-time availability tracking and seat management
    \item Database-backed user profiles and ride history
    \item Rating and reputation tracking with on-chain storage
    \item Emergency SOS alert system with Shamir's Secret Sharing
    \item Administrative dashboard for platform management
    \item Complaint management system with DPoS voting
    \item Eight deployed smart contracts for decentralized operations
    \item CP-ABE encryption service for vehicle information protection
    \item Merkle Tree audit anchoring for batch ride verification
    \item OpenStreetMap integration for location services
\end{itemize}

\subsection{Blockchain Security Features}

\begin{enumerate}
    \item \textbf{Soulbound Tokens (SBT)}: Non-transferable NFTs for academic identity
    \item \textbf{Zero-Knowledge Proofs (ZKP)}: Privacy-preserving driver verification
    \item \textbf{On-Chain Hash Integrity}: Tamper-proof ride agreement storage
    \item \textbf{DPoS Governance}: Decentralized delegate voting for disputes
    \item \textbf{Merkle Tree Audit Trail}: Efficient batch verification of ride history
    \item \textbf{CP-ABE Encryption}: Attribute-based access control for sensitive data
    \item \textbf{Shamir's Secret Sharing}: Threshold-based emergency data protection
\end{enumerate}

\section{Report Organization}

This report is organized as follows:

\begin{itemize}
    \item \textbf{Chapter 2: Literature Review} - Reviews existing blockchain applications in transportation and identifies gaps
    \item \textbf{Chapter 3: System Architecture} - Presents the three-tier architecture, technology stack, and database schema
    \item \textbf{Chapter 4: Implementation Details} - Describes user management, ride management, and safety features
    \item \textbf{Chapter 5: Blockchain Security Implementation} - Details all seven cryptographic security mechanisms
    \item \textbf{Chapter 6: Smart Contract Architecture} - Presents contract designs and deployment
    \item \textbf{Chapter 7: Conclusion} - Summarizes achievements and future directions
\end{itemize}
